\section{Experimental Method}
\label{sec:experiments}
\subsection{Populations and comparison conditions}
We analyze component measurements from the CAMP evaluation~\cite{campArtifact}
and new live transfer records~\cite{campTransferArtifact}. The memory-mapped
file access and feature-ablation studies use 33,516 tasks from GATK,
Minimap2, Sarek, Seqinspector, Taxprofiler, and Viralmetagenome. We use a
split of 28,490 development tasks and 5,026 evaluation tasks, using
chronological ordering within each workflow.

We first evaluate the contribution of memory-mapped file accesses by
comparing consumption measured through explicit reads alone with
consumption measured through explicit reads and memory-mapped accesses.
Both variants use the same evaluation tasks and measured memory peaks.
Each variant constructs its consumption history and trains its consumption
and memory models, with allocation policies selected independently on
development data. Table~\ref{tab:mmap} reports underallocation counts, total
requested memory, and median absolute percentage error (MdAPE). MdAPE
is the median absolute percentage error between predicted median peak
memory and measured peak memory.

\label{sec:training-calibration}
We then examine four feature variants: Access, which uses static
information and peak-memory history; Access with consumption history
$\mathbf{h}_{c,t}$; Access with predicted consumption $\hat c_t$; and CAMP,
which includes both consumption features. The retained LightGBM ensembles
and development-selected allocation policies are preserved.
Table~\ref{tab:features} reports underallocation counts, MdAPE, total
requested memory, and unused memory-time $W$ for these variants using
their recorded model configurations. The $W$ calculation uses tasks with
available runtimes.

The feature ablation runs on a 16-vCPU elastic Slurm worker, with 28~GB
of memory allocated to the job. The evaluation tasks, configured seeds,
and model hyperparameters remain unchanged from the original evaluation.
The reproduced Access and CAMP configurations yield 85 and 30
underallocations, respectively, compared with 87 and 31 reported in
WORKS~\cite{camp2026works}. We attribute these small differences to the
changed execution environment.

\begin{figure*}[!t]
\centering
\includegraphics[width=\textwidth]{\campReadyPath result_assets/audit_tradeoffs.pdf}
\caption{Two separate auditing studies. (a) Downstream underallocation-rate difference (selective minus random, in percentage points) over 17,067 tasks after the first update; negative values favor selective auditing. Bars are paired 95\% bootstrap confidence intervals. (b) Estimated auditing cost on 5,026 tasks with CAMP allocations fixed (5--50\% budgets shown); bars are 95\% bootstrap confidence intervals of five seed-level means. Panel (b) measures auditing overhead separately from the downstream replay in (a).}
\label{fig:audit}
\end{figure*}

\subsection{Downstream allocation and auditing cost}
Following RQ4 in WORKS~\cite{camp2026works}, we examine how selective
auditing affects subsequent memory allocations. Two identically
initialized CAMP predictors receive measurements selected through
selective or random auditing under equal budgets. Each predictor updates
between successive task batches. We compare underallocation counts and
total requested memory on tasks predicted after the first update, across
audit budgets of 5\%, 10\%, 20\%, 30\%, and 50\%, with five sampling seeds.
We retain 20\% as the primary downstream comparison.

A separate study estimates auditing cost while keeping the published CAMP
allocations fixed. At audit budget $B$, the fraction of tasks audited,
let $C_B$ denote the extra reserved memory-time caused by auditing, in
GiB-hours. For task $t$, $\hat y_t$ is CAMP's memory request in GiB;
$d_t^{\rm audit}$ and $d_t^{\rm base}$ are its runtimes in hours with and
without auditing. Summing over audited tasks ($q_t=1$) gives
\begin{equation}
 C_B=\sum_{t:q_t=1}\hat y_t\max(d_t^{\rm audit}-d_t^{\rm base},0).
\label{eq:audit_cost}
\end{equation}
The maximum counts only positive runtime overhead. Prediction, retraining,
and retry costs are excluded. Varying $B$ changes which tasks incur auditing
cost; the memory requests remain fixed. We evaluate costs at 5\%, 10\%,
15\%, 20\%, 25\%, 30\%, 50\%, and 100\% audit budgets. We report 95\% bootstrap confidence
intervals by resampling the five seed-level selective-minus-random
differences for downstream comparisons and the five seed-level cost
estimates for auditing cost.

\subsection{Live deployment and transfer protocol}
\label{sec:live-method}
The live study uses Bowtie2 with Nextflow and Slurm. Each stage executes
6,000 tasks across seven processes and produces 1,974 aligned chunks.
The original inputs are 38 paired-end sequencing runs from European
Nucleotide Archive (ENA) study PRJEB11501; the new inputs are 38 paired-end
sequencing runs from PRJNA486148. Selection uses metadata and compressed
size, without execution outcomes. The inputs total 15.14 and 15.20~GiB,
with no shared run IDs, sample IDs, file URLs, or FASTQ checksums. Read lengths nevertheless differ: 34 new runs
use 75-base reads and four use 102-base reads; original reads are 102 bases.

Nextflow configuration overlays enable the tracing described in
Section~\ref{sec:history}. Bowtie2's process-specific memory functions send
task metadata to CAMP's prediction service and return memory requests for
Slurm submission. Completion callbacks record task outcomes, and a
background learner publishes updated models for later predictions. This
integration preserves analytical tool commands while requiring
workflow-specific metadata bindings and deployment settings.

\emph{Warm-Transfer} executes cold, repeated-input warm, and new-input
stages in sequence. Cold begins with the seed-1996 offline model and
28,490 observations from other workflows, with no Bowtie2 history.
Each subsequent stage inherits the preceding terminal model, history,
and calibration. \emph{Direct-Transfer} starts the same new inputs from
that exact cold checkpoint, skipping the warm pass. Cold is not rerun.
The new-input branches therefore inherit 12,000 and 6,000 Bowtie2
observations, respectively. Each stage uses a fresh workflow directory;
completed tasks are not reused through Nextflow's cache.

A successful task outcome is eligible for model fitting when both
peak-memory and consumption measurements are available. The background
learner retrains after every 128 eligible outcomes. Completed-task records
enter CAMP's history
independently of these retraining updates. The latest 64 observations are
held out from model fitting and used to calibrate predictions from the
corresponding model version. A final update at the end of each stage
incorporates the remaining eligible outcomes. Failed attempts are retained
for diagnosis but excluded from model fitting. Both transfer branches
continue learning online and audit every task; the comparison therefore
evaluates the effect of an additional warm pass.

Three 16-vCPU, 64-GiB worker nodes expose 50~GiB each to Slurm. The control
node has 16 vCPUs and 64~GiB for Warm-Transfer, but 8 vCPUs and 32~GiB for
Direct-Transfer; the learner retains four CPUs. The software versions are
Nextflow 25.10.7, Bowtie2 2.4.2, and samtools 1.20. We calculate unused memory-time $W$ using CAMP's initial memory request
and the measured peak memory and runtime of the final successful attempt.
$W$ uses the measured runtime of successful attempts. Failed attempts
are excluded, and auditing, prediction, and retraining costs are not
separately accounted for.

There is one execution per branch and one new input study. Differences
in task-completion order can change the history available for subsequent
predictions and model updates.
